% Generated by scripts/build_reports.py; edit reviewed Markdown inputs.
% Compile with XeLaTeX.
% Options for packages loaded elsewhere
\PassOptionsToPackage{unicode}{hyperref}
\PassOptionsToPackage{hyphens}{url}
\PassOptionsToPackage{dvipsnames,svgnames,x11names}{xcolor}
\documentclass[
  11pt,
]{article}
\usepackage{xcolor}
\usepackage[margin=1in]{geometry}
\usepackage{amsmath,amssymb}
\setcounter{secnumdepth}{-\maxdimen} % remove section numbering
\usepackage{iftex}
\ifPDFTeX
  \usepackage[T1]{fontenc}
  \usepackage[utf8]{inputenc}
  \usepackage{textcomp} % provide euro and other symbols
\else % if luatex or xetex
  \usepackage{unicode-math} % this also loads fontspec
  \defaultfontfeatures{Scale=MatchLowercase}
  \defaultfontfeatures[\rmfamily]{Ligatures=TeX,Scale=1}
\fi
\usepackage{lmodern}
\ifPDFTeX\else
  % xetex/luatex font selection
\fi
% Use upquote if available, for straight quotes in verbatim environments
\IfFileExists{upquote.sty}{\usepackage{upquote}}{}
\IfFileExists{microtype.sty}{% use microtype if available
  \usepackage[]{microtype}
  \UseMicrotypeSet[protrusion]{basicmath} % disable protrusion for tt fonts
}{}
\makeatletter
\@ifundefined{KOMAClassName}{% if non-KOMA class
  \IfFileExists{parskip.sty}{%
    \usepackage{parskip}
  }{% else
    \setlength{\parindent}{0pt}
    \setlength{\parskip}{6pt plus 2pt minus 1pt}}
}{% if KOMA class
  \KOMAoptions{parskip=half}}
\makeatother
\usepackage{longtable,booktabs,array}
\usepackage{caption}
\captionsetup[table]{skip=6pt}
\newcounter{none} % for unnumbered tables
\usepackage{calc} % for calculating minipage widths
% Correct order of tables after \paragraph or \subparagraph
\usepackage{etoolbox}
\makeatletter
\patchcmd\longtable{\par}{\if@noskipsec\mbox{}\fi\par}{}{}
\makeatother
% Allow footnotes in longtable head/foot
\IfFileExists{footnotehyper.sty}{\usepackage{footnotehyper}}{\usepackage{footnote}}
\makesavenoteenv{longtable}
\setlength{\emergencystretch}{3em} % prevent overfull lines
\providecommand{\tightlist}{%
  \setlength{\itemsep}{0pt}\setlength{\parskip}{0pt}}
\usepackage{amsmath,amssymb}
\usepackage{xurl}
\usepackage{fancyhdr}
\pagestyle{fancy}
\fancyhf{}
\fancyhead[L]{\small Pure Thought | Revised research specifications}
\fancyfoot[L]{\small 2026-09-17}
\fancyfoot[R]{\small\thepage}
\setlength{\headheight}{15pt}
\setlength{\emergencystretch}{3em}
\allowdisplaybreaks
\usepackage{bookmark}
\IfFileExists{xurl.sty}{\usepackage{xurl}}{} % add URL line breaks if available
\urlstyle{same}
% fallback for those not using the hyperref driver hyperxmp:
\makeatletter
\@ifundefined{xmpquote}{\newcommand{\xmpquote}[1]{#1}}{}
\makeatother
\hypersetup{
  pdftitle={04. Lightcone bootstrap with controlled large-spin remainders},
  pdfauthor={Pure Thought research portfolio},
  colorlinks=true,
  linkcolor={blue},
  filecolor={Maroon},
  citecolor={Blue},
  urlcolor={blue},
  pdfcreator={LaTeX via pandoc}}

\title{04. Lightcone bootstrap with controlled large-spin remainders}
\author{Pure Thought research portfolio}
\date{Revised 17 September 2026}

\begin{document}
\maketitle
\thispagestyle{empty}
\clearpage

\section{04. Lightcone bootstrap with controlled large-spin remainders}

\textbf{Revised specification --- 2026-09-11.}
\href{https://github.com/rick-stevens-ai/Pure-Thought/blob/main/reviews/04-Modular-Lightcone-Bootstrap.md}{Detailed
critique} \ensuremath{\cdot}
\href{https://github.com/rick-stevens-ai/Pure-Thought/blob/main/archive/original-PRDs/04-Modular-Lightcone-Bootstrap.md}{Original
description} \ensuremath{\cdot}
\href{https://github.com/rick-stevens-ai/Pure-Thought/blob/main/reviews/EVIDENCE-STANDARD.md}{Evidence
standard}

\subsection{Plain-language guide}

\subsubsection{The problem in everyday terms}

A quantum correlation can be calculated by grouping the participating
excitations in different ways, and all valid groupings must agree. That
agreement can force some excitations to exist and restrict their
properties. This problem studies the pattern at large spin, where useful
approximations become available. The central challenge is to say not
only what the leading approximation predicts, but also how large its
error could be.

\subsubsection{Key terms}

\begin{itemize}
\item
  \textbf{Correlation function} --- A quantity describing how field
  measurements at different spacetime points are related.
\item
  \textbf{Four-point function} --- A correlation involving four field
  insertions; it can be reorganized into different pairings.
\item
  \textbf{Bootstrap} --- A method that constrains a theory by requiring
  its own consistency rules to agree, rather than starting from a
  complete microscopic model.
\item
  \textbf{Scaling dimension \ensuremath{\Delta}} --- A measure of how a
  field or excitation changes when all lengths are rescaled.
\item
  \textbf{Spin \ensuremath{\ell}} --- A label describing how an
  excitation transforms under rotations.
\item
  \textbf{Lightcone limit} --- A limit in which two spacetime points
  approach a separation that light could traverse; it organizes which
  contributions dominate the correlation function.
\item
  \textbf{Twist} --- The combination
  \ensuremath{\Delta}\ensuremath{-}\ensuremath{\ell}, useful for
  ordering contributions near a lightlike separation.
\item
  \textbf{Operator product expansion (OPE)} --- A rule expressing the
  effect of nearby field insertions as a sum of other operators with
  specified coefficients.
\item
  \textbf{Double-twist family} --- A sequence of large-spin operators
  whose dimensions approach a characteristic pattern related to two
  external operators.
\item
  \textbf{Remainder bound} --- A guaranteed limit on the terms omitted
  from an approximation.
\end{itemize}

\subsubsection{Why this matters}

Consistency can reveal features shared by many strongly interacting
theories, even when their detailed dynamics cannot be solved. Large-spin
patterns are also useful in theories proposed to describe gravity
through a lower-dimensional boundary. Controlled errors turn suggestive
asymptotic patterns into statements that can actually constrain a finite
calculation.

\subsubsection{What progress would mean}

The baseline would reproduce an established large-spin pattern with
unambiguous conventions. A stronger result would show how far that
pattern can deviate at specified spins, making later comparisons and
exclusions more trustworthy.

\subsection{Core question}

For one scalar four-point function in a unitary CFT with d\textgreater2,
what controlled bounds on large-spin double-twist data follow from
specified low-twist exchanges?

\subsection{Scope and assumptions}

Start in d=3 with an identical scalar \ensuremath{\phi} of fixed
dimension \ensuremath{\Delta}\ensuremath{\phi}. Specify the
normalization of C\_T, the stress-tensor OPE coefficient fixed by its
Ward identity, and an explicit twist gap above the low-twist exchanges
retained. Use a large-C\_T expansion only when a factorization
assumption and its error order are stated.

This challenge studies analytic large-spin behavior. Challenge 07
instead studies the spinning stress-tensor correlator. Torus modular
invariance is restricted to a separate d=2 variant.

\subsection{Mathematical target}

For a fixed scalar block convention, define
F\_\{\ensuremath{\Delta},\ensuremath{\ell}\}=v\textsuperscript{\{\ensuremath{\Delta}\ensuremath{\phi}\}g\_\{\ensuremath{\Delta},\ensuremath{\ell}\}(u,v)\ensuremath{-}u}\{\ensuremath{\Delta}\ensuremath{\phi}\}g\_\{\ensuremath{\Delta},\ensuremath{\ell}\}(v,u),
and use F\_id+\ensuremath{\Sigma}
\ensuremath{\lambda}\ensuremath{{}^{2}}F=0. Twist is
\ensuremath{\tau}=\ensuremath{\Delta}\ensuremath{-}\ensuremath{\ell}.
The double-twist families have
\ensuremath{\Delta}\_\{n,\ensuremath{\ell}\}=2\ensuremath{\Delta}\ensuremath{\phi}+2n+\ensuremath{\ell}+\ensuremath{\gamma}\_\{n,\ensuremath{\ell}\}
asymptotically.

Choose n=0 initially and derive a retained contribution
\ensuremath{\hat{\gamma}}\emph{\{0,\ensuremath{\ell}\}, together with a
bound
|\ensuremath{\gamma}}\{0,\ensuremath{\ell}\}\ensuremath{-}\ensuremath{\hat{\gamma}}\_\{0,\ensuremath{\ell}\}|\ensuremath{\le}R(\ensuremath{\ell})
for
\ensuremath{\ell}\ensuremath{\ge}\ensuremath{\ell}\ensuremath{{}_{0}}
under explicit spectral/Regge assumptions. If only an asymptotic series
is obtained, label it as such and do not manufacture a finite-spin
bound.

\subsection{Required outputs and proof obligations}

Export block normalizations, low-twist input, the derivation or
inversion formula, coefficients, and a decomposition of R into omitted
exchanges, expansion terms, and numerical errors. State clearly whether
the result is an identity, a conditional bound, or an asymptotic
estimate.

\subsection{Validation and rejection controls}

Use generalized-free-field crossing as an algebraic benchmark, while
noting that it is not by itself the target local CFT with a finite
stress tensor. Recover double-twist accumulation in a known analytic
example. Verify every uniform estimate on the specified cross-ratio
domain and all spins covered by the claim.

\subsection{Milestones and research extension}

Baseline: reproduce an established leading large-spin result with
conventions checked. Strong: obtain one explicit remainder estimate, or
identify the extra data needed to obtain it. Research extension: a
tighter conditional bound on low-twist exchanges or bulk couplings with
a stated holographic dictionary.

Milestones are acceptance gates, not calendar promises. Before a run,
freeze the concrete instance/domain, tolerance, compute budget, and
permitted dependencies. A bounded search that fails to find a result
must preserve an unresolved outcome.

\subsection{Starting instruction}

Work on the baseline above. First make every input and convention
explicit, then validate the reference case. Build the checker alongside
the calculation. Label numerical evidence, finite-model certificates,
and continuum theorems separately; return the strongest justified result
and identify the exact unmet proof obligation.

\subsection{Primary references and their role}

\begin{itemize}
\tightlist
\item
  \href{https://arxiv.org/abs/1212.3616}{Fitzpatrick et al., The
  Analytic Bootstrap and AdS Superhorizon Locality} --- Large-spin
  double-twist structure and analytic bootstrap.
\end{itemize}

References support definitions or limitations; the scope choices and
proposed milestones are editorial recommendations, not claims that these
results have already been obtained.

\subsection{Displayed formulation and claim boundary}

\href{https://github.com/rick-stevens-ai/Pure-Thought/blob/main/PRDs/04-Modular-Lightcone-Bootstrap.md}{Canonical
specification}.

\[F_{\Delta,\ell}=v^{\Delta_\phi}g_{\Delta,\ell}(u,v)
-u^{\Delta_\phi}g_{\Delta,\ell}(v,u),\qquad
F_{\mathrm{id}}+\sum_{\mathcal O}\lambda_{\phi\phi\mathcal O}^{2}F_{\Delta,\ell}=0.\]

\[\Delta_{n,\ell}=2\Delta_\phi+2n+\ell+\gamma_{n,\ell},\qquad
|\gamma_{0,\ell}-\widehat\gamma_{0,\ell}|\le R(\ell)\quad(\ell\ge\ell_0).\]

The inequality is a target requiring explicit spectral and Regge
assumptions and a proved remainder. An asymptotic series alone does not
imply this finite-spin inequality. Higher-dimensional crossing is not
torus modular invariance.

\subsection{Detailed review and rationale}

\href{https://github.com/rick-stevens-ai/Pure-Thought/blob/main/PRDs/04-Modular-Lightcone-Bootstrap.md}{Revised
problem} \ensuremath{\cdot}
\href{https://github.com/rick-stevens-ai/Pure-Thought/blob/main/archive/original-PRDs/04-Modular-Lightcone-Bootstrap.md}{Preserved
original} \ensuremath{\cdot}
\href{https://github.com/rick-stevens-ai/Pure-Thought/blob/f854d4e12b55c1f5b5c2a964b75f2d23496d981e/PRDs/04-Modular-Lightcone-Bootstrap.md}{Pinned
upstream source}

Line references below refer to the preserved original at commit
\texttt{f854d4e12b55}, not the rewritten file. ``Blocking'' means the
original issue can invalidate the scientific target or claimed
certificate; ``major'' means it materially changes scope, inference, or
acceptance criteria. These are critiques of specifications, not results
of running the embedded research-code sketches.

\subsubsection{Detailed findings}

\begin{enumerate}
\def\labelenumi{\arabic{enumi}.}
\item
  \textbf{Blocking --- modular invariance outside its domain (23, 45).}
  Thermal correlators on
  S\textasciicircum{}(d\ensuremath{-}1)\ensuremath{\times}S\ensuremath{{}^{1}}
  do not generally have the torus SL(2,Z) invariance asserted here. The
  revised d=3 problem uses crossing and Lorentzian analyticity only.
\item
  \textbf{Major --- spectrum terminology (16, 31--41).} A large
  holographic gap normally concerns a selected single-trace sector, not
  all nonconserved operators; double-twist operators remain.
  \ensuremath{\Delta}\_gap\ensuremath{-}J\_max is not a definition of a
  minimum twist. The scalar unitarity lower bound is also not a
  holographic gap prediction.
\item
  \textbf{Major --- asymptotics versus finite bounds.} Lightcone
  expansions alone do not supply certified finite-spin exclusion plots.
  A remainder or additional input is required.
\item
  \textbf{Major --- overlap and overclaim (29, 53--60).} The original
  mixes analytic scalar bootstrap, spinning bootstrap, modular
  constraints, and uniqueness of Einstein gravity. The replacement
  isolates a calculable question and preserves gravity interpretations
  as conditional extensions.
\end{enumerate}

\subsubsection{What the revision preserves and changes}

The central motivation is retained, but the executable question is now:
\textbf{For one scalar four-point function in a unitary CFT with
d\textgreater2, what controlled bounds on large-spin double-twist data
follow from specified low-twist exchanges?} The revised scope chooses a
concrete starting model; those choices are proposed research-design
decisions, not facts inferred from the original. The baseline can
succeed without resolving the full open research program.

The replacement also removes the original implementation sketches from
the active specification. They remain in the archive for comparison, but
have not been made executable or certified. New implementations should
derive their equations from the revised definitions rather than copy
unchecked prototypes.

\subsubsection{Acceptance criteria to use instead}

Baseline: reproduce an established leading large-spin result with
conventions checked. Strong: obtain one explicit remainder estimate, or
identify the extra data needed to obtain it. Research extension: a
tighter conditional bound on low-twist exchanges or bulk couplings with
a stated holographic dictionary.

\subsubsection{Sources supporting the corrections}

\begin{itemize}
\tightlist
\item
  \href{https://arxiv.org/abs/1212.3616}{Fitzpatrick et al., The
  Analytic Bootstrap and AdS Superhorizon Locality} --- Large-spin
  double-twist structure and analytic bootstrap.
\end{itemize}

\section{Evidence and execution standard}

This standard applies to the revised specifications. It is a proposed
contract for future research runs; this editorial revision does not
supply the scientific certificates themselves.

\subsection{What ``pure thought'' means here}

Use explicitly stated mathematical models and deductions without fitting
conclusions to experimental outcomes. Synthetic examples, exact model
inputs, analytic references, numerical search, and established software
are allowed. Physical constants, electronic-structure approximations,
constitutive laws, and empirical parameter sets must be identified when
used. A model-derived result is conditional on that model.

The original ban on legacy software conflicts with its proposed reliance
on NumPy, SciPy, CVXPY, nauty, and other libraries. Reimplementing every
optimizer and algebra system is not necessary for scientific
independence. Prefer a small independently checkable certificate and a
documented trusted base. A fresh implementation can be a project
objective, but must not be confused with a mathematical proof or novelty
claim.

\subsection{Four evidence levels}

{\def\LTcaptype{none} % do not increment counter
\begin{longtable}[]{@{}
  >{\raggedright\arraybackslash}p{(\linewidth - 4\tabcolsep) * \real{0.3333}}
  >{\raggedright\arraybackslash}p{(\linewidth - 4\tabcolsep) * \real{0.3333}}
  >{\raggedright\arraybackslash}p{(\linewidth - 4\tabcolsep) * \real{0.3333}}@{}}
\toprule\noalign{}
\begin{minipage}[b]{\linewidth}\raggedright
Label
\end{minipage} & \begin{minipage}[b]{\linewidth}\raggedright
What it establishes
\end{minipage} & \begin{minipage}[b]{\linewidth}\raggedright
What it does not establish
\end{minipage} \\
\midrule\noalign{}
\endhead
\bottomrule\noalign{}
\endlastfoot
Numerical evidence & A reproducible computation at stated precision,
sample size, and discretization & Exact truth, continuum coverage, or a
universal theorem \\
Finite-model certificate & An exact/validated statement about the
specified finite model or finite relaxation & An omitted tail,
thermodynamic limit, or physical realization \\
Conditional mathematical theorem & A proof over the stated domain with
every hypothesis and error controlled & Truth outside those assumptions
or a demonstrated physical system \\
Formal proof & A checked theorem in a named proof assistant, with
imported axioms/dependencies listed & Correct physical modeling merely
by virtue of formalization \\
\end{longtable}
}

Formalization is an optional extra deliverable unless the specific
project makes it mandatory. High precision, JSON output, solver status
and small residuals do not by themselves confer any of these proof
levels.

\subsection{Claim directions}

For minimization over a larger feasible set, the optimum is a lower
bound on the original minimum. A physical feasible point supplies an
upper bound. For maximization the directions reverse. Track feasibility
of the primal/dual witnesses separately from their objective values.

Infeasibility of an inner ansatz excludes only that ansatz.
Infeasibility of an outer relaxation can exclude the original problem if
the map into the relaxation is proved and all approximations preserve
necessity. Passing necessary tests means ``not excluded''; existence
needs a constructive sufficient argument.

A smallest observed example is not a minimality theorem. A complete
catalog must define the universe, equivalence relation, boundaries, and
proof that pruning misses nothing. A finite-size numerical trend does
not prove an asymptotic limit.

\subsection{Minimum run manifest}

Before calculation, record the problem ID and revision hash, exact
inputs and conventions, quantified domain, equivalence relation, target
observable, accepted evidence level, tolerances, and resource limits.
Record dependencies and precision/rounding mode. Where randomization is
used, preserve seeds, trial counts and raw sufficient statistics.
Resource limits are planning decisions, not evidence that a mathematical
solution cannot exist.

An output bundle should include:

\begin{itemize}
\tightlist
\item
  The statement being claimed, including all assumptions and
  quantifiers.
\item
  Construction or proof object, plus an independently runnable verifier
  and its trust assumptions.
\item
  Numerical/error ledger separating roundoff, discretization,
  truncation, optimization, sampling and modeling error where relevant.
\item
  Positive reference cases and negative controls that the verifier must
  reject.
\item
  Coverage: included domain, omitted domain, unresolved regions and
  termination reason.
\end{itemize}

These are requirements for future implementations, not a promise of
polynomial-size certificates for every problem.

\subsection{Checking specific certificate types}

\textbf{LP/SDP/SOS:} verify equality constraints, objective direction,
and cone feasibility using exact arithmetic or outward-rounded
intervals. Rationalizing a numerical solution can break PSD constraints.
Correct residuals with justified norm bounds; do not silently project
and reuse the old objective. Distinguish a proof of a polynomial
approximation from a proof for the original function.

\textbf{Interval numerics:} state function enclosures, domains and
outward rounding. A small root box needs existence and, if claimed,
uniqueness. Cover the complement for a global enumeration. Include
singularities and tails rather than discarding them numerically.

\textbf{SAT/SMT:} a solver's UNSAT label is not a certificate. Preserve
a proof supported by a checker and the proof that the encoding, symmetry
reduction and blocked solutions match the mathematical universe. DRAT
applies to an appropriate Boolean CNF proof, not automatically to every
real-arithmetic or SDP claim.

\textbf{Statistics:} specify the estimand, dependence assumptions,
confidence procedure and bias. Zero observed failures does not prove
zero failure probability. MCMC error bars based on estimated
autocorrelation are not generally rigorous mixing guarantees.

\subsection{Milestones and novelty}

Baseline means a working, independently validated result on a concrete
reference problem. Strong means a broader or sharper result with its
claimed evidence level. Research extension means a possible novel result
whose feasibility and significance need review. A negative theorem can
be valuable; an unsuccessful search can also be useful if its scope is
documented, but the two are not interchangeable.

Do not promise ``publication-quality'' in a fixed number of months.
Before claiming novelty, conduct a focused current literature comparison
against the exact theorem/model/metric. The sources in this revision
check key definitions and limitations; they are not an exhaustive
novelty audit across all 30 fields.

\end{document}
